% Section content template for individual Educative lessons
% This file contains the actual content for a specific section
% Generated from Educative API section content

% Section: Quiz: Design Patterns
% Section ID: 5218544979804160
% Section Slug: quiz-design-patterns
% Generated: 2025-12-11T06:16:39.708479

\subsection{MCQs}\label{MCQs}
\\
Challenge yourself by solving the following quiz questions.

\subsection*{Design Patterns}
\vspace{0.3cm}
\noindent
\textbf{Question 1:} \textbf{(Select all that apply.)}
\\
Which statement is \emph{true} about design patterns?
\\
\textit{(Select all that apply.)}
\\[0.2cm]
\begin{itemize}
 \item \textbf{[CORRECT]} Design patterns provide one of the best solutions, as they have been derived and optimized over time by various experienced programmers.
 \item \textbf{[CORRECT]} Design patterns provide a clean and elegant solution to a large problem by avoiding repetition in code.
 \item \textbf{[CORRECT]} Design patterns can complicate the application's architecture if they are poorly managed.
 \item It is necessary to use design patterns in the code even if a simpler approach exists.
\end{itemize}
\vspace{0.3cm}
\noindent
\textbf{Question 2:} What best describes the design pattern types most commonly used in software products?
\\[0.2cm]
\begin{itemize}
 \item Functional, Relational, Sequential
 \item \textbf{[CORRECT]} Creational, Structural, Behavioral
 \item Abstract, Concrete, Static
 \item Object, Method, Class
\end{itemize}
\vspace{0.3cm}
\noindent
\textbf{Question 3:} Suppose you are building a notification system that can send emails, SMS, or push notifications based on user settings. Which pattern would you use to create the appropriate notification object at runtime, and why?
\\[0.2cm]
\begin{itemize}
 \item Singleton to ensure that only one notification service exists.
 \item \textbf{[CORRECT]} Factory to create different types of notification objects as needed.
\\[0.1cm]
\textit{Explanation:} 
The Factory pattern provides a way to create objects without specifying the exact class. When the type of object needed varies at runtime (email, SMS, push), a Factory can instantiate the right object based on input or configuration, making the system flexible and easily extensible.
 \item Adapter to make notification classes compatible.
 \item Composite to combine all notification methods.
\end{itemize}
\vspace{0.3cm}
\noindent
\textbf{Question 4:} How does the Singleton pattern differ from other creational patterns?
\\[0.2cm]
\begin{itemize}
 \item It always creates a new object.
 \item \textbf{[CORRECT]} It restricts a class to only one instance, providing a global access point.
\\[0.1cm]
\textit{Explanation:} 
The Singleton pattern ensures that a class has only one instance and provides a single global point of access to that instance. This is useful for resources like configuration managers, loggers, or database connections, where having multiple instances could cause problems.
 \item It allows creating families of related objects.
 \item It decorates objects with new functionality.
\end{itemize}
\vspace{0.3cm}
\noindent
\textbf{Question 5:} You are developing a tool that builds complex HTML documents with optional components (like headers, footers, and multiple sections). Which creational pattern would make your code more flexible and manageable?
\\[0.2cm]
\begin{itemize}
 \item The Adapter pattern
 \item The Decorator pattern
 \item The Facade pattern
 \item \textbf{[CORRECT]} The Builder pattern
\\[0.1cm]
\textit{Explanation:} 
The Builder pattern is used to construct complex objects step by step, especially when different representations or configurations are needed. It separates the construction logic from the object's representation, making the code easier to maintain and extend when optional or complex components are involved.
\end{itemize}
\vspace{0.3cm}
\noindent
\textbf{Question 6:} You have a \texttt{Coffee} class. You want to allow adding extra ingredients (like milk, sugar, or whipped cream) dynamically, without changing the base \texttt{Coffee} class code. Which structural pattern should you use, and why?
\\[0.2cm]
\begin{itemize}
 \item Proxy to control access to the \texttt{Coffee} object.
 \item \textbf{[CORRECT]} Decorator to add new features to \texttt{Coffee} objects dynamically.
\\[0.1cm]
\textit{Explanation:} 
The Decorator pattern allows you to add new behaviors or features to objects at runtime, by ``wrapping'' them with additional functionality. This avoids modifying the original class, making it ideal for adding optional features like coffee toppings.
 \item Composite to treat individuals and groups of \texttt{Coffee} objects the same.
 \item Adapter to convert \texttt{Coffee} to another interface.
\end{itemize}
\vspace{0.3cm}
\noindent
\textbf{Question 7:} What best describes the Adapter pattern?
\\[0.2cm]
\begin{itemize}
 \item It provides a way to create families of related objects.
 \item It defines a one-to-many dependency.
 \item \textbf{[CORRECT]} It allows incompatible interfaces to work together by translating one interface into another.
\\[0.1cm]
\textit{Explanation:} 
The Adapter pattern acts as a bridge between two incompatible interfaces, allowing existing code to work with new code without modification. It 's commonly used to make old code compatible with new APIs or third-party libraries.
 \item It restricts object creation to a single instance.
\end{itemize}
\vspace{0.3cm}
\noindent
\textbf{Question 8:} You are building a blog platform. When a new article is published, all subscribers should be notified automatically (by email, SMS, etc.), without the blog code having to know the details of each notification method. Which pattern should you use?
\\[0.2cm]
\begin{itemize}
 \item \textbf{[CORRECT]} The Observer pattern
\\[0.1cm]
\textit{Explanation:} 
The Observer pattern lets you define a subscription mechanism so that multiple objects (observers) can be notified automatically of changes or events in another object (the subject), decoupling the publisher from the subscribers.
 \item The Singleton pattern
 \item The Proxy pattern
 \item The Command pattern
\end{itemize}
\vspace{0.3cm}
\noindent
\textbf{Question 9:} What problem does the Command pattern solve, and how is it typically implemented?
\\[0.2cm]
\begin{itemize}
 \item It restricts object instantiation; uses a static method.
 \item \textbf{[CORRECT]} It encapsulates a request as an object, allowing parameterization and queuing of requests.
\\[0.1cm]
\textit{Explanation:} 
The Command pattern turns requests or actions into standalone objects, which can then be queued, logged, or undone. This decouples the sender from the receiver and adds flexibility for features like undo/redo, macros, and delayed execution.
 \item It allows objects to be decorated; uses wrapping classes.
 \item It allows objects to be iterated; uses \texttt{next}/\texttt{hasNext} methods.
\end{itemize}
\vspace{0.3cm}
\noindent
\textbf{Question 10:} How is data synchronization between the UI and the underlying data model typically achieved in the MVVM (Model-View-ViewModel) architectural pattern?
\\[0.2cm]
\begin{itemize}
 \item Through manual updates in the Controller
 \item By direct method calls from the View to the Model
 \item \textbf{[CORRECT]} By using data binding, which keeps the View and ViewModel in sync automatically
\\[0.1cm]
\textit{Explanation:} 
In MVVM, data binding connects the View and the ViewModel so changes in the ViewModel are instantly reflected in the UI, and vice versa. This automatic synchronization is a key feature, especially in modern frameworks that support declarative UI.
 \item By letting the Presenter dictate all UI changes
\end{itemize}



% End of section content